
\subsubsection{3.1 Binomial Variance Analysis}

Consider a problem with T test cases and a candidate solution generated by a model. Let q $\in$ [0,1] denote the per-test-case success probability, treated as independent across tests for a given solution.

\textbf{Ratio reward.} The ratio reward is:

$$R_{\text{ratio}} = \frac{\text{tests\_passed}}{T} \in [0, 1]$$

If tests\_passed ~ Binomial(T, q), then the expected within-group variance of a single rollout's ratio reward is:

$$\mathbb{E}[\text{Var}(R_{\text{ratio}})] = \frac{q(1-q)}{T}$$

This is maximized at q = 0.5, where E[Var] = 1/(4T), and remains above 0.2/T for q $\in$ [0.2, 0.8].

\textbf{Binary reward.} The binary reward is:

$$R_{\text{binary}} = \mathbf{1}[\text{all } T \text{ tests pass}] \in \{0, 1\}$$

The probability of passing all T tests is $q^T$, so the expected within-group variance is:

$$\mathbb{E}[\text{Var}(R_{\text{binary}})] = q^T(1 - q^T)$$

This collapses rapidly with T because the probability of simultaneous success across all T tests decreases exponentially.

\textbf{Variance ratio.} Define the variance ratio:

$$\rho(q, T) = \frac{\mathbb{E}[\text{Var}(R_{\text{ratio}})]}{\mathbb{E}[\text{Var}(R_{\text{binary}})]} = \frac{q(1-q)}{T \cdot q^T(1-q^T)}$$

For representative parameter values: at T = 5 and q = 0.5, $\rho \approx$ 1.65. At T = 5 and q = 0.3, $\rho \approx$ 17.3. The ratio grows substantially with T and peaks at lower q values within the partial-tractability window. For q $\in$ [0.3, 0.55] and T = 5, $\rho$ ranges from approximately 5 to 20. This derivation follows from standard Binomial distribution properties; the contribution is its application to GRPO's group-relative advantage normalization.

\textbf{Significance for GRPO.} In GRPO, advantages are $A_i$ = (r\_i - mean(r)) / std(r). When std(r) $\rightarrow$ 0, all advantages collapse to zero regardless of reward values. Higher within-group reward variance directly produces larger non-degenerate advantage magnitudes, which in turn produce larger gradient updates. The variance analysis above quantifies the mechanism by which ratio rewards are expected to produce more informative GRPO gradients than binary rewards on problems in the partial-tractability regime.

Figure 1 displays the analytical variance comparison.

\begin{figure}[htbp]
\centering
\includegraphics[width=\columnwidth]{figures/fig_variance_advantage.png}
\caption{Figure 1: Theoretical variance advantage of R\_ratio over R\_binary}
\label{fig:fig_variance_advantage}
\end{figure}

\textit{Figure 1. Expected within-group variance E[Var(r)] as a function of problem tractability q for R\_ratio (solid) and R\_binary (dashed) across T $\in$ {3, 5, 8, 10} test cases (left), and variance ratio $\rho$(q, T) (right). The prescreening target window S\_term $\in$ [0.3, 0.55] is shaded. In this region, R\_ratio achieves a 5--$20\times$ variance advantage over R\_binary for T = 5.}

\subsubsection{3.2 Problem Tractability Formalization}

\textbf{Definition (Group Tractability Score).} For a problem p and k sampled solutions, the group tractability score is:

$$S_{\text{term}}(p) = \frac{1}{k} \sum_{i=1}^{k} \mathbf{1}[\text{tests\_passed}(s_i, p) \geq 1]$$

S\_term measures the fraction of rollouts that pass at least one test case. It is a group-level empirical estimate of the probability that a randomly sampled solution achieves at least partial credit.

\textbf{Target window.} S\_term $\in$ [0.3, 0.55]. The lower bound of 0.3 ensures that at least 2--3 rollouts per group of k=8 produce non-zero rewards, maintaining within-group reward heterogeneity. The upper bound of 0.55 excludes problems that are near-trivially solvable by the model, where binary reward variance is already non-negligible. Problems with S\_term = 0 (no rollout passes any test) are explicitly excluded: under either reward function, all rewards are 0 and no gradient flows.

\textbf{Analytical basis for the window.} The variance ratio $\rho$(q, T) is maximized at lower values of q within the tractable range. For T = 5, $\rho$ reaches its peak advantage in the range q $\in$ [0.3, 0.4]. The upper bound of 0.55 ensures the window excludes the regime where binary rewards are already informative. The window parameters are analytically motivated; empirical calibration after SFT checkpoint generation is recommended to verify their optimality.

\subsubsection{3.3 Prescreening Protocol}

The prescreening protocol selects training problems for which the current model has partial tractability before any GRPO gradient updates are applied.

\textbf{Algorithm 1: Pass@k Prescreening for GRPO Problem Selection}

\begin{verbatim}
Input:  Problem set P, model M
        k=8, τ=0.8, max_new_tokens=1024, seed=42
        S_low=0.3, S_high=0.55, T_min=3

Output: Tractable subset P* ⊆ P

1. Filter P: retain p with T(p) ≥ T_min
2. For each problem p in P:
   a. Sample k solutions from M at temperature τ
   b. Execute each solution against test cases (5s timeout per test)
   c. Compute S_term(p) = (1/k) Σ_i 1[tests_passed(s_i, p) ≥ 1]
3. P* = {p : S_low ≤ S_term(p) ≤ S_high}
4. Return P*
\end{verbatim}

The group size k=8 follows the standard GRPO group size from \citep{Shao2024}. The temperature $\tau =0.8$ balances exploration with code coherence. The T\_min=3 filter ensures at least three distinct test outcomes are possible under R\_ratio, providing meaningful reward granularity.

\subsubsection{3.4 Reward Functions}

\textbf{Ratio reward:}
$$R_{\text{ratio}}(s, p) = \frac{\sum_{j=1}^{T(p)} \mathbf{1}[\text{test}_j(s) = \text{PASS}]}{T(p)} \in [0, 1]$$

\textbf{Binary reward:}
$$R_{\text{binary}}(s, p) = \mathbf{1}\!\left[\sum_{j=1}^{T(p)} \mathbf{1}[\text{test}_j(s) = \text{PASS}] = T(p)\right] \in \{0, 1\}$$

\textbf{GRPO advantage computation:} $A_i$ = (r\_i $- \mu _{r})$ / ($\sigma _r + \epsilon )$. When $\sigma _r$ = 0, advantages are undefined; in practice a small $\epsilon$ is added. When $\sigma _r \approx$ 0, advantages are near-zero for all rollouts and no meaningful learning occurs. Prescreening aims to ensure $\sigma _r$ remains bounded away from zero by selecting problems where reward variance is structurally non-zero.

\subsubsection{3.5 Experimental Pipeline}

The experimental pipeline comprises two stages.

\textbf{Stage 1 --- Prescreening (h-e1):} Apply Algorithm 1 to the problem set; compute S\_term for each problem; filter to the target window [0.3, 0.55]; evaluate pre-registered gate metrics.

\textbf{Stage 2 --- GRPO Training Comparison (h-m1 through h-m4):} On the prescreened subset, run parallel GRPO training with R\_ratio versus R\_binary; record advantage distributions, gradient covariance, gradient SNR, and Zero Reward Fraction (ZRF) curves. Stage 2 is gated on Stage 1 passing.

Figure 2 shows the pipeline architecture.

\begin{figure}[htbp]
\centering
\includegraphics[width=\columnwidth]{figures/fig_prescreening_pipeline.png}
\caption{Figure 2: Prescreening pipeline architecture}
\label{fig:fig_prescreening_pipeline}
\end{figure}

\textit{Figure 2. Full prescreening pipeline. APPS introductory problems are loaded (difficulty=0, $T\geq 3)$, k=8 rollouts are generated per problem, S\_term is computed, and problems with S\_term $\in$ [0.3, 0.55] are retained for GRPO training. The gate evaluates whether fraction\_k\_pass\_ge1, pct\_groups\_above\_1.5x, and n\_prescreened meet pre-registered thresholds.}

